\scanpage{35}
\[
\underset{\unsure{\chi|?|\mathfrak g}}{\text{\unsure{Algèbres|Algèbre|Alg Spéc}}}\ (\mathcal M)
\ \unsure{\subseteq|\simeq}\mathbb P(\widetilde{\unsure{M|J|\mathfrak g}})
\simeq \mathbb P^\vee(M)\simeq
\bigl[\check P=\mathrm{Dr}(P)\bigr]\simeq \operatorname{Div}^{+}_{2}(X)
\]
% Note below the second projective-space expression, joined to it by an inclusion sign.
\[
\rotatebox{90}{$\supset$}
\]
\begin{center}
Ss-gpes lisses\\
connexes \unsure{comm.|commut.}\\
\deleted{\unsure{de}} de dim 1\\
dans $G$\\[2pt]
$\unsure{\operatorname{Jac}|J^{\cdot\cdot}|\mathcal T\in\cdots}(G)$
\end{center}

$\operatorname{Cent}(J)=J$ -- le centralisateur de \deleted{\unsure{tte}} \unsure{tte|toute} section de $J$ telle que $g^2\ne1$ en chaque fibre est $J$.

$U=X-Y$ torseur sous $J$

% slanted
\margin{\deleted{\illegible{1}} \unsure{si|en} \unsure{caractér.} \unsure{ou|en} car 2, $\unsure{J_s\simeq\mathbf G_{a,s}|J_s\simeq\mathrm{Ga}_s}$}

\deleted{Dans} $N(J)/J\subset\underline{\operatorname{Aut}}(J)$ \deleted{il y a} la section $\sigma=-\mathrm{id}_J$ est dans $N(J)/J$

\deleted{Soit} \deleted{La} l'image inverse $\unsure{\bar J|\bar{\jmath}|\overline{\mathcal T}}$ de $\sigma$ dans $N(J)$ \deleted{\unsure{est un}} \inserted{qui est un} \deleted{$J$-torseur} $J$-torseur \deleted{$X-Y$, en est une} est comme $X-Y$. \deleted{\unsure{On a}} \deleted{\unsure{c}} \deleted{section ce qui}
\[
\begin{aligned}
(X-Y)&\longrightarrow\bar J\\
J&\xrightarrow{2\,\mathrm{id}_J} J
\end{aligned}
\]
% archive sheet number: 34
